\section{Methodology}\label{sec:model}

\subsection{Competitive Environment}\label{sec:model-payoff}

We adapt the multi-contest framework of
\citet{xu_zenou_zhou_conflict_2022} to firms allocating capacity
across activities serving overlapping customer demands.
Capacity generates appeal at each demand node, and a firm's appeal relative
to rivals and outside alternatives determines its share of the available
contribution margin.
In the static game, each firm chooses capacity to maximize that margin net
of capacity costs, taking rivals' allocations as given.

Let firm $i=1,\ldots,N$ assign capacity $q_{ia}\geq0$ to activity
$a=1,\ldots,m$, located at $x_a\in\mathbb R^d$.
Customer demand arises at $L$ nodes $z_\ell\in\mathbb R^d$.
These locations describe economic activities and demands, while article-vector
coordinates will represent their reported profiles in
Section~\ref{sec:adjustment-design}.
To distinguish scale from business mix, we represent the allocation by a
measure $Q_i$ with total mass $S_i$ and normalized composition $P_i$:
\begin{equation}\label{eq:position}
  Q_i=\sum_a q_{ia}\delta_{x_a},\qquad
  S_i=\sum_{a}q_{ia},\qquad
  P_i=\sum_a p_{ia}\delta_{x_a},\quad p_{ia}=q_{ia}/S_i,
\end{equation}
where $\delta_{x_a}$ is a unit point mass at $x_a$ and $P_i$ is defined when
$S_i>0$.
The vector $p_i$ records the composition of simultaneous activities, not
probabilities of selecting a single activity.
Keeping total capacity $S_i$ separate allows the same disturbance to change
operating scale and business mix differently.
Capacities are in common model units, and the joint allocation
$q=(q_1,\ldots,q_N)$ lies in $\mathcal Q=\mathbb R_+^{Nm}$
(Appendix~\ref{app:proofs}).

Capacity in one activity can serve several related customer segments.
The coefficient $K_{\ell a}=\exp\{-\tau_A\|x_a-z_\ell\|^2/d\}$ measures how
effectively activity $a$ serves demand node $\ell$, declining with the
distance between them, and a higher $\tau_A$ makes it decline faster and narrows
the range of demands a given activity serves.
A firm's appeal at a node is the sum of its activity capacities weighted by
that effectiveness, $A_i=Kq_i\in\mathbb R_+^L$, and each firm faces the
combined appeal of its rivals and of an outside alternative, measured in the
same units:
\begin{equation}\label{eq:appeal}
  A_{i\ell}(q_i)=\sum_{a}K_{\ell a}q_{ia},\qquad
  B_{i\ell}(q_{-i})=b_\ell+\sum_{j\ne i}A_{j\ell}(q_j),\qquad b_\ell>0.
\end{equation}
The outside appeal $b_\ell$ represents alternatives beyond the modeled firms
and keeps the share denominator positive even when no modeled firm
serves a node.
Own capacity increases a firm's access to a demand opportunity, while rivals
divide that opportunity among more participants.

At each demand node, firms compete for a contribution-margin opportunity
$r_{\ell t}\geq0$ available at date $t$, valued in cash-flow units.
Exogenous free cash flow shocks change these available contribution margins.
Firms' initial exposures and subsequent capacity choices determine the
resulting changes in their cash flows.
A firm receives the fraction of that opportunity given by its share of total
appeal, including the outside alternative, so its operating payoff is
contribution margin less the persistent cost
$C_i$ of maintaining capacity:
\begin{equation}\label{eq:firm-problem}
  \pi_i(q_i,q_{-i};r_t)=
  \underbrace{\sum_\ell r_{\ell t}
  \frac{A_{i\ell}(q_i)}{A_{i\ell}(q_i)+B_{i\ell}(q_{-i})}}
  _{\text{contribution margin from current opportunities}}
  -\underbrace{C_i(q_i)}_{\text{cost of positioning there}}.
\end{equation}
Each firm chooses nonnegative capacity to maximize this stylized unlevered
operating cash flow, taking its rivals' allocations as given.
Its contribution-margin component, denoted $\mathcal M_i$
(Appendix~\ref{app:payoff}), adapts the proportional-share contest of
\citet{xu_zenou_zhou_conflict_2022}: capacities generate appeal through $K$,
and $b_\ell$ represents alternatives outside the modeled firms.
Appendix~\ref{app:contest-inheritance} gives the adaptation and its required
assumptions.
Doubling an opportunity weight doubles its contribution to the margin at an
unchanged profile, whereas doubling every firm's capacity need not double the
margin because shares approach saturation.
Summed over firms, the margins at a node never exceed $r_{\ell t}$: firms
divide a finite pool, and crowding works through forgone contribution margin.

The proportional-appeal share has a discrete-choice interpretation when
capacity counts varieties or outlets and customers value nearby activities.
Appendix~\ref{app:logit} derives that interpretation following
\citet{train_discrete_choice_2009}, with opportunity value
combining market size and fixed per-sale margin.

\subsection{Costs and Equilibrium}\label{sec:model-equilibrium}

Firms differ in the activities they can expand cheaply.
A capability center $h_i\in\mathbb R^d$, a location in the activity space,
makes nearby activities less costly at the margin,
with the penalty for distance governed by $\gamma\geq0$.
The persistent capability cost is
\begin{equation}\label{eq:cost}
  C_i(q_i)=c_i^\top q_i+\tfrac12q_i^\top Hq_i,
  \quad c_{ia}=c_0+\gamma\|x_a-h_i\|^2/d,
  \quad H=\eta K^c+\nu I,
\end{equation}
where $c_0,\eta\geq0$, $\nu>0$, and $K^c_{ab}=\exp\{-\tau_c\|x_a-x_b\|^2/d\}$.
Dividing squared distances by $d$ expresses both distance penalties per
activity-space coordinate.
The linear term represents persistent comparative advantage, the Gaussian
congestion term makes simultaneous expansion in related activities costly,
and $\nu$ raises the marginal cost of concentrating capacity at any one site.
Increasing $\gamma$ makes departures from a firm's capability center more
expensive, so capability centers give firms different costs of responding to
the same opportunity.
These are costs of operating the allocation, including when it is unchanged.

Together, the proportional-appeal margin and the cost specification define
the operating payoff in \eqref{eq:firm-problem}.
The common Gaussian matrix in \eqref{eq:cost} is the numerical specialization
used below.
For the equilibrium argument, we allow a symmetric firm-specific matrix $H_i$
subject to the following conditions.

\begin{assumption}[Demand access and expansion costs]\label{ass:primitives}
  Outside appeal satisfies $b_\ell>0$, opportunity values and appeal
  coefficients
  satisfy $r_\ell\geq0$ and $K_{\ell a}\geq0$, and linear costs
  satisfy $c_{ia}\geq0$.
  Each symmetric cost matrix obeys $H_i\succeq\nu I$ for a common $\nu>0$.
\end{assumption}

The Gaussian congestion matrix is positive semidefinite, so the numerical
specification satisfies the curvature condition.
At equilibrium, marginal contribution equals marginal capacity cost on
every positive-capacity activity, while entry into an unused activity
cannot raise payoff.
To establish that these joint conditions determine a unique allocation,
Appendix~\ref{app:equilibrium} adapts the contest and variational-inequality
methods of
\citet{xu_zenou_zhou_conflict_2022} and
\citet{parise_ozdaglar_network_games_2019}.

\begin{proposition}[Unique equilibrium and the symmetry
  boundary]\label{prop:equilibrium}
  Under Assumption~\ref{ass:primitives}, the proportional-share capacity game
  has a unique Nash equilibrium on $\mathcal Q$.
  Identical firm primitives imply identical equilibrium capacity vectors.
  The independent-firm economy, obtained by setting $B_{i\ell}=b_\ell$,
  also has a unique equilibrium, consisting of the firms' individual optima.
\end{proposition}

Positive outside appeal keeps demand shares defined, while strictly curved
costs discipline scale and ensure uniqueness
(Appendix~\ref{app:equilibrium}).
Uniqueness also explains the symmetry boundary: crowding can spread aggregate
activity across opportunities, but firms with identical primitives spread it
identically.
Capability differences supply the heterogeneous incentives that separate
firms' business mixes in the illustrative economies.

\subsection{Sequential Adjustment and Valuation}\label{sec:model-dynamics}

Starting from the static equilibrium, firms incur additional costs when
changing their inherited capacity allocations.
In each period they choose capacities simultaneously, with adjustment costs
relative to the preceding allocation $q_{i,t-1}$.
The stage choice is
\begin{equation}\label{eq:anchored-choice}
  q_{it}\in\arg\max_{z_i\geq0}\left\{
  \pi_i(z_i,q_{-i,t};r_t)-\frac{\kappa}{2}\|z_i-q_{i,t-1}\|^2\right\},\qquad\kappa\geq0.
\end{equation}
Adjustment costs make operating responses gradual, as in the costly
repositioning problem studied by \citet{sweeting_positioning_2013}.
The quadratic form makes every stage a strongly monotone Nash problem
(Appendix~\ref{app:anchored-stage}).
Firms maximize current operating payoff net of adjustment expenditure,
without a continuation-value term.
Financial prices reflect the resulting choices and do not feed back into
investment.

At an interior initial equilibrium $\bar q$, let $J$ be the joint Jacobian
of marginal capacity costs minus marginal contributions, and let $\mathcal G$
map opportunity changes into changes in marginal incentives.
A local capacity deviation $x_t$ generated by opportunity deviation
$s_t v$ obeys
\begin{equation}\label{eq:anchored-response}
  (J+\kappa I)x_t=\kappa x_{t-1}+\mathcal G v s_t.
\end{equation}
Adjustment costs make earlier investment decisions persist: a firm responds
to new opportunities from the capacity allocation it has already built.
Because $J$ includes rivals' effects on marginal returns, competitors'
simultaneous adjustments also shape that response.
At a boundary equilibrium, equation~\eqref{eq:anchored-response} applies
locally while the active set remains unchanged.
The nonlinear simulation solves the actual level game in
\eqref{eq:anchored-choice} and checks its departure from this affine
response.

We value the cash flows generated by this adjustment path using a constant
discount factor $\delta\in(0,1)$.
With baseline payout $\bar\pi_i>0$, the ex-payout price is
$\bar V_i=\delta\bar\pi_i/(1-\delta)$.
For a unit permanent opportunity innovation, let $x_0=0$ and let $x_h$ for
$h\geq1$ solve \eqref{eq:anchored-response} with $s_h=1$.
The unexpected total-return loading is
\begin{equation}\label{eq:same-path-value}
  g_i=\frac{1}{\bar V_i}\sum_{h=0}^{\infty}\delta^h
  \left\{\partial_r\pi_i(\bar q;\bar r)v+D_q\pi_i(\bar q;\bar r)x_h\right\}.
\end{equation}
The first term values the opportunity at existing capacities; the second
values the cash-flow consequences of subsequent operating responses.
Investors can incorporate this anticipated path when the news arrives,
before capacity reaches its new allocation.
Adjustment expenditure has zero first derivative at the unchanged origin,
so it adds no first-order payout term to this valuation.

The discounted capacity sum solves
$[J+(1-\delta)\kappa I]\sum_h\delta^h x_h
=\delta\mathcal Gv/(1-\delta)$, valuing the full response without truncating
its horizon.
In the spirit of \citet{berk_green_naik_growth_options_1999}, operating
opportunities enter firm valuations through the cash flows they generate,
which here follow the supplied sequence of Nash choices under
constant discounting.

\subsection{Reporting and the Directional Prediction}\label{sec:model-observation}

Financial news records firms' activities through the coverage they receive.
Investors seeking to value future cash flows have reason to follow material
investment, product and technology decisions, while publishers and corporate
disclosures influence which developments become observable
\citep{gentzkow_media_bias_2006}.
To connect this reporting process to operating decisions, we assume that an
activity's share of news is proportional to its capacity weighted by its
reporting intensity.
Capacity measures the activity available to be reported, while reporting intensity
allows coverage per unit of capacity to differ across activities and time.
For fixed component probability laws $\nu_a$, this gives
\begin{equation}\label{eq:reporting-law}
  C_{i,t}=\sum_a\omega_{ia,t}\nu_a,\qquad
  \omega_{ia,t}=\frac{\lambda_{ia,t}q_{ia,t}}
  {\sum_b\lambda_{ib,t}q_{ib,t}},\qquad\lambda_{ia,t}>0.
\end{equation}
Common intensity across activities gives $\omega_{ia,t}=p_{ia,t}$, so the
article law follows the firm's business mix under equal relative attention.
Changing relative attention changes reported composition even at fixed
capacities, whereas reciprocal changes in capacity and intensity can
preserve the same observed law.
Text-based product measures \citep{hoberg_phillips_2016} provide a precedent
for studying activities through text, while the reporting weights specify
the maintained connection to operating capacity here.
The annual observation pools the dated article laws before evaluating a
transport distance.

In the restricted two-opportunity economy below, an improvement in one
opportunity relative to the other locally raises each firm's capacity in the
improving activity and lowers its capacity in the other.
Every firm therefore increases its share of the improving opportunity,
although the joint response includes the capacity choices of its rivals.
Firms with greater initial exposure to the improving opportunity have larger
return responses to the shock.
Under common reporting across activities, a firm locally increasing its share
of the improving opportunity approaches the origin profiles of firms with
larger shares of that opportunity and recedes from those with smaller shares.

Consider diversified firms with distinct interior shares $p_i$ in that
opportunity, common total capacity $Q$ and $\sum_i p_i=N/2$.
With identity activity loading, common baseline opportunity value $v$ and
denominator $D=b+NQ/2>2Q$, each node's response Jacobian is
$J=d_\perp I+u\mathbf1^\top$, where
$d_\perp=\nu+v/D^2$ and $u_i=v(D-2q_i)/D^3>0$.
The anchored response has two persistence modes, $\kappa/(d_\perp+\kappa)$
and $\kappa/(d_A+\kappa)$ with $d_A=d_\perp+\sum_i u_i$.

Valuing the cash flows along this adjustment path gives the return response
$g(p)=G(2p-1)/[1+(2p-1)^2]$ for $G>0$.
The return response $g(p)$ increases with the firm's initial exposure $p$.

For the unit-axis component laws
$\nu_-=(\delta_{e_1}+\delta_{-e_1})/2$ and
$\nu_+=(\delta_{e_2}+\delta_{-e_2})/2$, every mixture has zero centroid, yet
\begin{equation}\label{eq:axis-observation}
  W_2^2(C(p),C(q))=2|p-q|.
\end{equation}
The target therefore approaches the initial news distributions of reference
firms with larger return responses to the shock.
Under the timing and observation conditions below, this implies a nonnegative
expected local covariance between approach and relative returns.
Write $A_{ijy}$ for the local decrease in population Wasserstein squared
distance of target $i$ toward
reference $j$ over year $y$, and $R_{ijy}$ for the corresponding article-weighted
excess-return measure, with observed counterparts defined in
Section~\ref{sec:adjustment-design}.

\begin{proposition}[Directional movement--return
  alignment]\label{prop:directional-alignment}
  Maintain the interior two-opportunity construction of
  Appendix~\ref{app:canonical-mechanism} and the unit-axis observation laws
  under common reporting across activities.
  Hold origin profiles and article timing fixed, with nonnegative source-year
  timing weights and no earlier-window articles strictly after a current-year
  source.
  Suppose primitive source-date innovations are centered, independent across
  dates and have finite variances, and maintain the zero cross-moment
  restrictions on financial proxy error, initial gaps and observation error.
  Then, for each target $i$ and destination year $y$,
  \[
    E[\operatorname{Cov}_{j\ne i}(R_{ijy},A_{ijy})]\geq0.
  \]
  Fixed nonnegative pooling of these moments preserves the sign.
\end{proposition}

Appendix~\ref{app:canonical-mechanism} proves the result by deriving the
ordered share and valuation responses from the same joint operating path.
News may be shared across firms: independence concerns source dates, while
firms' different exposures generate their relative responses.
